\documentclass[../main.tex]{subfiles}
\ifSubfilesClassLoaded{\setcounter{chapter}{1}}{}
\begin{document}

\chapter{Lingkungan Python 3, Editor, dan Identifier}

\begin{subcpmk}
  \item Sub-CPMK 1.1: Menginstal interpreter Python 3 dan menjalankan skrip dari editor atau lingkungan Jupyter sesuai petunjuk praktikum
\end{subcpmk}

\SesiRPS{2}{1.1}{$\sim$170 menit}{CPMK-1}

\section{Sarana dan prasarana}

Komputer laboratorium dengan Python 3.10+ (atau versi yang ditetapkan prodi), terminal, editor/IDE, Jupyter Notebook (opsional).

\section{Prosedur kerja}

\begin{enumerate}[leftmargin=*]
  \item Verifikasi versi interpreter (\texttt{python --version}).
  \item Ikuti contoh REPL, skrip \texttt{.py}, dan Jupyter pada modul ini.
  \item Kerjakan aktivitas dan unggah artefak sesuai konvensi Bab 0 / Lampiran F.
\end{enumerate}

\section{Pengenalan Python 3}

Python adalah bahasa tingkat tinggi yang banyak dipakai untuk pembelajaran, sains data, otomasi, dan pengembangan web. Versi yang digunakan pada modul ini adalah \textbf{Python 3}. Pastikan perintah \texttt{python --version} atau \texttt{python3 --version} menampilkan 3.x.

\section{Menjalankan Program}

\subsection{REPL (Read--Eval--Print Loop)}

Buka terminal, jalankan \texttt{python} atau \texttt{python3}, lalu ketik ekspresi:

\begin{lstlisting}[caption=Contoh sesi REPL]
>>> 2 + 3
5
>>> print("Halo laboratorium")
Halo laboratorium
\end{lstlisting}

\subsection{Skrip Berkas \texttt{.py}}

Simpan kode ke berkas \texttt{halo.py}, lalu jalankan dari folder yang sama:

\begin{lstlisting}[caption=Skrip sederhana]
# halo.py
print("Halo dari berkas")
\end{lstlisting}

Perintah: \texttt{python halo.py} (atau \texttt{python3 halo.py}).

\subsection{Jupyter Notebook}

Jupyter menyediakan sel kode dan dokumentasi.\newline Setelah instalasi (\texttt{pip install notebook}),\newline jalankan \texttt{jupyter notebook} dan buat buku kerja baru. Sel kode dieksekusi dengan Shift+Enter.

\section{Editor dan Konvensi}

Gunakan editor yang mendukung indentasi konsisten (VS Code, PyCharm, Neovim, dll.). Python memakai \textbf{indentasi} untuk blok, bukan kurung kurawal.

\section{Identifier dan Gaya Penamaan}

Nama variabel/fungsi (\textit{identifier}) memakai huruf, angka, dan garis bawah; tidak boleh diawali angka. Konvensi umum: \texttt{snake\_case} untuk fungsi dan variabel, hindari kata cadangan seperti \keyword{if}, \keyword{for}, \keyword{class}.

\begin{contoh}
Nama yang baik: \texttt{nilai\_tugas}, \texttt{hitung\_rata}. Hindari: \texttt{1nilai}, \texttt{if} (konflik kata kunci).
\end{contoh}

\section{Studi Kasus dan Contoh Kode}

Berikut adalah contoh-contoh program Python yang mengimplementasikan konsep-konsep pada modul ini.

\subsection{Kode: \texttt{hello.py}}
\begin{lstlisting}[language=Python,caption={\texttt{hello.py}}]
"""Menampilkan Hello dunia (cuplikan minimal)."""


def main() -> None:
    print("Hello, dunia")
    input("Tekan Enter untuk melanjutkan...")


if __name__ == "__main__":
    main()
\end{lstlisting}

\subsection{Kode: \texttt{hello2.py}}
\begin{lstlisting}[language=Python,caption={\texttt{hello2.py}}]
"""Hello dengan variabel string dan integer (setara hello2.pas)."""


def main() -> None:
    # Nama Program : hello2.pas - cuplikan dari modul Pascal
    shello = "Hello, dunia"
    i = 1234
    print(shello)
    print("Nilai i = ", i)
    input("Tekan Enter untuk melanjutkan...")


if __name__ == "__main__":
    main()
\end{lstlisting}

\section{Referensi sesi}

\begin{itemize}[leftmargin=*]
  \item \textbf{[8, 9]} \textit{The Python Tutorial} (Python 3, EN/ID).
  \item \textbf{[10]} Python Wiki, \textit{BeginnersGuide}.
  \item \textbf{[11]} Severance, C. \textit{Python for Everybody (PY4E)}.
\end{itemize}

\begin{aktivitas}
  \item Verifikasi versi Python di stasiun lab dan catat pada lembar observasi.
  \item Buat berkas \texttt{perkenalan.py} yang mencetak nama NIM dan kelas Anda.
  \item Buka Jupyter, buat satu sel yang mencetak hasil \texttt{7 * 6}, dan simpan buku kerja sebagai \texttt{latihan01.ipynb}.
  \item Diskusi singkat: apa perbedaan menjalankan kode di REPL versus berkas?
\end{aktivitas}

\begin{latihan}
  \item Jelaskan dua cara menjalankan program Python yang telah Anda praktikkan.
  \item Sebutkan aturan penamaan identifier yang valid dan beri masing-masing satu contoh salah.
  \item Apa fungsi komentar dengan tanda \texttt{\#}? Tulis contoh satu baris komentar pada kode.
  \item Refleksi: bagian mana dari lingkungan lab yang paling membingungkan pada pertemuan ini?
\end{latihan}

\begin{asesmen}
\textbf{Instrumen singkat (formatif)}

\begin{enumerate}
  \item Perintah mana yang menampilkan versi Python di terminal?
  \begin{enumerate}
    \item \texttt{python --help}
    \item \texttt{python --version}
    \item \texttt{pip install python}
    \item \texttt{jupyter --run}
  \end{enumerate}
  \item Manakah identifier yang \textbf{tidak} valid?
  \begin{enumerate}
    \item \texttt{data2024}
    \item \texttt{\_nilai}
    \item \texttt{for\_loop}
    \item \texttt{2nilai}
  \end{enumerate}
\end{enumerate}

\textbf{Rubrik}: jawaban benar + bukti pengumpulan berkas \texttt{.py} atau \texttt{.ipynb} sesuai instruksi dosen.
\end{asesmen}

\begin{checklist}
  \item Saya dapat menjalankan Python di terminal atau Jupyter
  \item Saya dapat membuat dan menjalankan berkas \texttt{.py}
  \item Saya memahami aturan identifier dasar
  \item Saya menyimpan pekerjaan sesuai format penamaan yang diminta
\end{checklist}

\begin{rangkuman}
Bab ini membahas instalasi/verifikasi Python 3, menjalankan kode lewat REPL, skrip, dan Jupyter, serta konvensi identifier.

\textbf{Poin Kunci:} Python 3; jalankan skrip dengan \texttt{python nama\_berkas.py}; indentasi penting; nama variabel mengikuti aturan identifier.

\textbf{Kata kunci:} \code{REPL}, \code{Jupyter}, \code{.py}, \code{identifier}
\end{rangkuman}

\end{document}
